\chapter{AnalysisDirector — Multi-Phase Orchestrator (Phase D)}
\label{ch:analysis_director}

% =============================================================================
%  Chapter 62: AnalysisDirector, state threading, phase gating,
%              solve_incremental state synchronisation, and test
%              infrastructure improvements
% =============================================================================

\section{Motivation}

Real structural analysis workflows rarely consist of a single solver
invocation.  A typical seismic assessment, for example, involves:

\begin{enumerate}
  \item \textbf{Static preload:} gravity and imposed dead loads via
        \texttt{NonlinearAnalysis} ($p: 0 \to 1$).
  \item \textbf{Dynamic excitation:} earthquake ground motion via
        \texttt{DynamicAnalysis} ($t: 0 \to T$).
  \item \textbf{Static redistribution:} post-event settling via
        \texttt{NonlinearAnalysis} again.
\end{enumerate}

Each phase may use a different solver engine and different configurations,
yet they share structural state: the displacement at the end of the
preload becomes the initial condition for the dynamic phase, and the
post-earthquake displacement feeds the redistribution phase.

Chapters~59--61 built the individual components (incremental control,
step director, steppable-solver concept, analysis state transfer).
Phase~D assembles them into a single orchestrator:
\texttt{AnalysisDirector}.


\section{Design Decisions}

\subsection{Type-Erased Phase Callbacks}

The director is \emph{not} templated on solver types.  Instead, each
phase is a \texttt{std::function<AnalysisState(const AnalysisState\&)>}:
the callback receives the output of the previous phase (or a
default-constructed state for the first phase) and returns its own
output.

This design gives maximum flexibility:
\begin{itemize}
  \item The caller decides how to create, configure, and destroy each
        solver.
  \item Different phases can use entirely different engines
        (\texttt{NonlinearAnalysis}, \texttt{DynamicAnalysis}, or any
        future solver satisfying \texttt{SteppableSolver}).
  \item Boundary conditions, observers, and loads can be freely modified
        between phases without the director knowing their details.
\end{itemize}

The trade-off is that the director cannot impose compile-time guarantees
on solver types.  This is acceptable because the director is a
\emph{workflow} coordinator, not a solver abstraction.

\subsection{Borrowed State and Lifetime}

\texttt{AnalysisState} stores \emph{borrowed} PETSc \texttt{Vec}
references (see Chapter~61).  This implies a critical lifetime
constraint:

\begin{quote}
\textbf{Rule:} The solver that produced an \texttt{AnalysisState} must
remain alive while any subsequent phase reads its state.
\end{quote}

In practice, this means solvers that participate in state transfer
should be constructed \emph{outside} the phase callbacks, so their
\texttt{Vec} resources outlive the inter-phase handoff.  The tests
in this chapter demonstrate this pattern explicitly.

\subsection{Failure Semantics}

A phase signals failure by returning an \texttt{AnalysisState} with
\texttt{displacement == nullptr}.  The director supports two modes:

\begin{itemize}
  \item \textbf{Abort-on-failure (default):} the phase sequence stops
        immediately; remaining phases are not executed.
  \item \textbf{Continue-on-failure:} all phases execute regardless of
        prior failures.  Useful for robustness studies or post-mortem
        analysis.
\end{itemize}

\subsection{Phase Gating}

An optional \texttt{PhaseGateCallback} can be attached to any phase.
The gate receives the \emph{previous} phase's \texttt{PhaseOutcome} and
returns a boolean: \texttt{true} to proceed, \texttt{false} to skip.

Typical gates include:
\begin{itemize}
  \item Only run a dynamic phase if the preload reached full load:
        \texttt{prev.state.time >= 0.9}.
  \item Only run a redistribution phase if the excitation succeeded:
        \texttt{prev.status == PhaseStatus::Succeeded}.
\end{itemize}


\section{Architecture}

\subsection{Core Types}

\begin{verbatim}
namespace fall_n {

enum class PhaseStatus { Succeeded, Failed, Skipped };

struct PhaseOutcome {
    std::string    name;
    PhaseStatus    status{PhaseStatus::Skipped};
    double         elapsed_seconds{0.0};
    AnalysisState  state{};
};

struct DirectorReport {
    std::vector<PhaseOutcome> phases;
    bool        all_succeeded() const noexcept;
    std::size_t num_succeeded() const noexcept;
    std::size_t num_failed()    const noexcept;
    void        print_summary() const;
};

using PhaseCallback     = std::function<AnalysisState(const AnalysisState&)>;
using PhaseGateCallback = std::function<bool(const PhaseOutcome&)>;

} // namespace fall_n
\end{verbatim}

\subsection{AnalysisDirector API}

\begin{verbatim}
class AnalysisDirector {
public:
    // Registration
    AnalysisDirector& add_phase(std::string name, PhaseCallback cb);
    AnalysisDirector& add_phase(std::string name, PhaseCallback cb,
                                PhaseGateCallback gate);
    AnalysisDirector& set_continue_on_failure(bool);

    // Execution
    DirectorReport run();

    // Queries
    std::size_t num_phases() const noexcept;
};
\end{verbatim}

The \texttt{add\_phase()} overloads return \texttt{*this} for chaining.

\subsection{Execution Algorithm}

\texttt{run()} iterates over registered phases in order:

\begin{enumerate}
  \item \textbf{Gate check:} if the phase has a gate and the previous
        phase's outcome exists, evaluate the gate.  If it returns
        \texttt{false}, record the phase as \texttt{Skipped} and
        continue to the next phase.
  \item \textbf{Execute:} call the phase callback with the current
        \texttt{AnalysisState}.  Measure wall-clock time via
        \texttt{std::chrono::steady\_clock}.
  \item \textbf{Classify:} if the returned state has a non-null
        displacement, mark \texttt{Succeeded} and update the current
        state.  Otherwise, mark \texttt{Failed}.
  \item \textbf{Abort check:} if the phase failed and
        \texttt{continue\_on\_failure} is \texttt{false}, break.
\end{enumerate}

The initial state for the first phase is a default-constructed
\texttt{AnalysisState} (all null, time~0, step~0).


\section{Bug Fix: \texttt{solve\_incremental} State Synchronisation}
\label{sec:solve_incremental_sync}

During Phase~D testing, a latent inconsistency was discovered:
\texttt{NonlinearAnalysis::solve\_incremental()} uses a \emph{local}
variable \texttt{p\_done} for bookkeeping, but never updates the
\emph{member} \texttt{p\_done\_} used by \texttt{get\_analysis\_state()}.

This meant that calling \texttt{get\_analysis\_state()} after
\texttt{solve\_incremental()} would report $p = 0$ and step $= 0$
regardless of the actual terminal state.

\subsection{Fix Applied}

After the solve loop and element state update, the batch path now
synchronises with the single-step state:

\begin{verbatim}
model_->update_elements_state();

// Synchronise single-step state so get_analysis_state()
// reflects the terminal position even when called after
// the batch solve_incremental() path.
p_done_     = p_done;
step_count_ = steps_done;
\end{verbatim}

This ensures that \texttt{get\_analysis\_state()} is consistent
regardless of whether the solve was performed via the batch API
(\texttt{solve\_incremental}) or the step-by-step API
(\texttt{begin\_incremental + step()}).


\section{Usage Example}

\begin{verbatim}
// Setup model (shared across phases)
Model<Policy> model(domain, material);
model.fix_x(0.0);
model.setup();
model.set_density(2400.0);  // concrete

// Create solvers outside callbacks for lifetime safety
NonlinearAnalysis<Policy> nl{&model};
DynamicAnalysis<Policy>   dyn{&model};

AnalysisDirector director;

director.add_phase("gravity", [&](const AnalysisState&) {
    nl.solve_incremental(10);
    return nl.get_analysis_state();
});

director.add_phase("earthquake", [&](const AnalysisState& prev) {
    dyn.set_initial_displacement(prev.displacement);
    dyn.set_time_step(0.005);
    dyn.step_to(30.0);
    return dyn.get_analysis_state();
});

auto report = director.run();
report.print_summary();
\end{verbatim}

Output:
\begin{verbatim}
AnalysisDirector — Phase Summary
┌────┬────────────────────────────────┬───────────┬────────────┐
│  # │ Phase                          │ Status    │ Time (s)   │
├────┼────────────────────────────────┼───────────┼────────────┤
│  1 │ gravity                        │ OK        │      0.042 │
│  2 │ earthquake                     │ OK        │     12.340 │
└────┴────────────────────────────────┴───────────┴────────────┘
All 2 phases completed successfully.
\end{verbatim}


\section{Verification}

\begin{center}
\begin{tabular}{clr}
\hline
\textbf{\#} & \textbf{Test} & \textbf{Assertions} \\
\hline
1  & Empty director                                        & 5  \\
2  & Single phase (static)                                 & 7  \\
3  & Two-phase state threading                             & 6  \\
4  & Static preload $\to$ Dynamic excitation (multi-engine)& 9  \\
5  & Abort-on-failure                                      & 4  \\
6  & Continue-on-failure mode                              & 6  \\
7  & Phase gate callback (pass and block)                  & 5  \\
8  & DirectorReport summary queries                        & 4  \\
9  & Three-phase pipeline (preload $\to$ excite $\to$ redistribute) & 7  \\
10 & Phase timing (\texttt{elapsed\_seconds > 0})          & 1  \\
\hline
   & \textbf{Total}                                        & \textbf{54} \\
\hline
\end{tabular}
\end{center}

\textbf{Test~4} is the integration highlight: it performs a static
preload with \texttt{NonlinearAnalysis}, transfers the displacement
via \texttt{AnalysisState} to \texttt{DynamicAnalysis}, and verifies
that the dynamic solver correctly inherits the deformed configuration.

\textbf{Test~9} demonstrates the full three-phase pipeline
(preload $\to$ excitation $\to$ redistribution) on a shared model,
confirming that state threads correctly through all transitions.

All 52/52 suite tests pass.


\section{Test Infrastructure Improvements}

During Phase~D, four pre-existing test failures (from Phase~C) were
resolved.  All four had the same root cause: \textbf{missing data files}.

\subsection{MSHReader Defensive Fix}

The mesh reader \texttt{MSHReader} (in \texttt{src/mesh/gmsh/ReadGmsh.hh})
silently printed an error to \texttt{std::cerr} when a mesh file was not
found, then proceeded to return to the caller with an uninitialised stream.
Downstream code (in \texttt{GmshDomainBuilder}) would then dereference
null pointers and crash with a segmentation fault.

\textbf{Fix:} replaced the \texttt{std::cerr} message with a
\texttt{throw std::runtime\_error(...)}, providing a clear diagnostic
and preventing silent null-pointer dereferences.

\subsection{Mesh File Generation}

The test data files (\texttt{.msh} meshes and earthquake records) are
excluded from version control via \texttt{.gitignore}.  A Python script
\texttt{scripts/generate\_test\_data.py} was created to regenerate them
using the Gmsh Python API (v4.15.1):

\begin{description}
  \item[\texttt{tests/validation\_cube.msh}] \hfill \\
    Second-order hexahedral mesh ($2 \times 2 \times 2$, 8 Hex27 elements,
    125 nodes).  Generated from \texttt{validation\_cube.geo} via
    \texttt{gmsh.model.mesh.generate(3)} + \texttt{setOrder(2)}.

  \item[\texttt{data/input/Beam\_LagCell\_Ord1\_30x6x3.msh}] \hfill \\
    First-order hexahedral mesh ($30 \times 6 \times 3$, 540 Hex8 elements,
    868 nodes).  Box geometry $[0,10] \times [0,0.4] \times [0,0.8]$,
    transfinite meshing via Gmsh API.

  \item[\texttt{data/input/earthquakes/el\_centro\_1940\_ns.dat}] \hfill \\
    Synthetic earthquake record (501 points, $\Delta t = 0.02$\,s,
    PGA $= 0.3194g$ at $t = 4.80$\,s).  Jennings envelope with random-phase
    sinusoids mimicking the El Centro 1940 NS component.
\end{description}

\textbf{Key insight:} hand-writing MSH 4.1 format files
programmatically appeared to produce correct meshes (node/element counts
matched) but caused \texttt{SNES DIVERGED\_FNORM\_NAN} at iteration~0.
Using the Gmsh API guarantees correct node ordering and element
connectivity, resolving the issue completely.


\section{Files}

\begin{description}
  \item[\texttt{src/analysis/AnalysisDirector.hh}] \hfill \\
    \texttt{PhaseStatus}, \texttt{PhaseOutcome}, \texttt{DirectorReport},
    \texttt{PhaseCallback}, \texttt{PhaseGateCallback},
    \texttt{AnalysisDirector} class.

  \item[\texttt{src/analysis/NLAnalysis.hh}] \hfill \\
    Modified: \texttt{solve\_incremental()} now synchronises
    \texttt{p\_done\_} and \texttt{step\_count\_} at completion.

  \item[\texttt{src/analysis/Analysis.hh}] \hfill \\
    Added: \texttt{AnalysisDirector.hh} include.

  \item[\texttt{src/mesh/gmsh/ReadGmsh.hh}] \hfill \\
    Modified: \texttt{std::cerr} $\to$ \texttt{throw std::runtime\_error}
    on file-not-found.

  \item[\texttt{tests/test\_analysis\_director.cpp}] \hfill \\
    10 tests, 54 assertions.

  \item[\texttt{scripts/generate\_test\_data.py}] \hfill \\
    Gmsh API mesh generator + synthetic earthquake record generator.

  \item[\texttt{CMakeLists.txt}] \hfill \\
    Registered \texttt{fall\_n\_analysis\_director\_test}.
\end{description}


\section{Phase Summary}

\begin{center}
\begin{tabular}{clcl}
\hline
\textbf{Phase} & \textbf{Component} & \textbf{Chapter} & \textbf{Status} \\
\hline
A & Single-step API (\texttt{DynamicAnalysis})   & 60 & Complete \\
B & StepDirector + factories                     & 60 & Complete \\
C & SteppableSolver concept + NL step API        & 61 & Complete \\
D & AnalysisDirector orchestrator                & 62 & Complete \\
\hline
\end{tabular}
\end{center}

The four-phase architecture provides a complete workflow construction
kit: individual steps (Phase~A), conditional breakpoints (Phase~B), a
unified solver concept (Phase~C), and a multi-phase orchestrator
(Phase~D).  Together, they enable arbitrary analysis sequences — from a
simple static solve to a multi-phase seismic assessment — with full
state transfer and lifecycle control.
